\chapter{Preliminaries}
\label{chap:preliminaries}

This chapter introduces the fundamental concepts and definitions required to understand the proposed methodology. 
It defines the problems of image super-resolution and image segmentation, summarizes essential deep learning concepts, 
and describes the objective functions (loss function) and evaluation metrics used throughout this thesis.

%-----------------------------------------------------------
\section{Image Super-Resolution}
\label{sec:prelim_sr}

Image super-resolution (SR) aims to reconstruct a high-resolution (HR) image from a given low-resolution (LR) input. 
Traditional interpolation methods such as bicubic upsampling often yield blurry and over-smoothed results, 
as they rely on fixed mathematical assumptions rather than learned patterns. 
Deep learning models learn data-driven mappings that reconstruct missing high-frequency details from low-resolution inputs.
The SR task is inherently ill-posed, since multiple HR images can correspond to the same LR observation. 
Therefore, the goal of modern SR models is to learn a mapping function 
\[
Y = f_\theta(X),
\]
where \(X\) is the LR input, \(Y\) the ground-truth HR image, and \(f_\theta\) a neural network with learnable parameters \(\theta\) 
that minimizes the difference between predicted and target images.
\begin{figure}[t]
    \centering
    \includegraphics[width=\textwidth]{LaTeX_ThesisEnglish/Images/super_resolution.png}
    \caption{Qualitative super-resolution example. 
    From left to right: ground-truth high-resolution (HR) image, low-resolution (LR) image upsampled by interpolation, 
    model prediction, absolute difference between HR and prediction, and edge-difference (Sobel) map. 
    The bottom row shows zoomed-in crops highlighting the differences between HR, LR-upsampled, and the predicted image.}
    \label{fig:sr-qualitative-car}
\end{figure}



%-----------------------------------------------------------
\section{Image Segmentation}
\label{sec:prelim_segmentation}

Image segmentation assigns a class label to every pixel in an image, dividing it into semantically meaningful regions. 
In biomedical imaging, segmentation is used to delineate anatomical or pathological structures such as lesions or organs. 
In this thesis, segmentation serves primarily as an auxiliary task to evaluate the cross-domain generalization of the proposed adaptive-depth U-Net. 
Because both super-resolution and segmentation employ encoder–decoder architectures with skip connections, a single unified framework can effectively support both \cite{ronneberger2015unet,long2015fcn}.

\subsection{Skin Lesions in Dermoscopic Imaging}
Skin lesions are localized abnormalities of the skin that can indicate benign or malignant conditions. 
In dermoscopic images, the lesions exhibit large variations in color, texture, and sharpness of the boundaries, often visually overlapping with the surrounding tissue. 
Accurate segmentation of the lesion area provides a useful benchmark for model robustness and spatial precision, but is not the primary focus of this work. 
The ISIC-2017 dataset offers standardized dermoscopic images and corresponding lesion masks for objective evaluation of segmentation performance \cite{phan2023skin}.

%-----------------------------------------------------------
\section{Core Deep Learning Components}
\label{sec:prelim_cnn}

This section introduces the core building blocks used in the U-Net architecture.

\subsection{Convolutional Neural Networks (CNNs)}

\begin{itemize}
    \item \textbf{Convolutional Layer:} 
    The core operation of CNNs applies a set of learnable filters (kernels) across the input image to produce feature maps. 
    This enables the network to capture hierarchical spatial features such as edges, textures, and shapes.
    
    \item \textbf{Activation Functions:} Nonlinear activation functions allow networks to model complex relationships. 
    The most common choice, and the one used in this thesis, is the Rectified Linear Unit (ReLU), defined as \( f(x) = \max(0, x) \).
    
    \item \textbf{Downsampling:} 
    Pooling or strided convolutions are conventionally used to reduce spatial dimensions, increase the receptive field and enable extraction of contextual information of higher level.
    
    \item \textbf{Upsampling:} 
    Transposed convolutions or interpolation layers increase spatial resolution to reconstruct fine details of the image during the decoding stage.
\end{itemize}

\subsection{The U-Net Architecture}

The U-Net architecture, introduced by Ronneberger et al.\ (2015), was originally developed for biomedical image segmentation \cite{ronneberger2015unet}. 
It follows a symmetric encoder–decoder structure:

\begin{itemize}
    \item \textbf{Encoder (Contracting Path):} 
    A stack of convolutional and pooling layers that progressively capture global context.
    
    \item \textbf{Decoder (Expanding Path):} 
    A series of upsampling and convolutional layers that reconstruct spatial resolution and localize details.
    
    \item \textbf{Skip Connections:} 
    Feature maps from each encoder level are concatenated with their corresponding decoder level. 
    These skip connections preserve high-frequency information and help the network recover spatial precision lost during downsampling.
\end{itemize}

This architecture is effective for both super-resolution and segmentation, 
as it combines semantic context with pixel-level detail reconstruction \cite{ronneberger2015unet}.

\subsection{Residual Learning}

Residual learning, introduced in ResNet \cite{he2016resnet}, simplifies the training of deep networks by letting layers learn residual corrections instead of full transformations. 
Rather than directly predicting an output \( H(x) \), the model learns a residual function \( F(x) = H(x) - x \), 
and reconstructs the final output as \( H(x) = F(x) + x \). 
This approach improves gradient flow and makes optimization more stable.

In super-resolution, the same principle is applied by predicting only the missing high-frequency details instead of the full high-resolution image. 
The residual head in this thesis implements this through the \texttt{ClipAdd} layer:
\[
\hat{Y} = \text{ClipAdd}(X_{LR}, \text{Residual}(X_{LR})),
\]
which adds the predicted residual to the low-resolution input and clips the output to ensure all pixel values remain within the valid intensity range \([0,1]\).

\subsection{Interpolation-Based Scaling Layers}
\label{sec:resize_layers}

Traditional U-Net architectures rely on fixed pooling and upsampling operations that are discrete and non-differentiable with respect to scale. 
Inspired by interpolation-based pooling concepts proposed by Wang~\cite{wang2020interpolation}, this thesis uses differentiable resizing layers that allow smooth and flexible scaling of feature maps.

\begin{figure}[t]
    \centering
    \includegraphics[width=\textwidth]{LaTeX_ThesisEnglish/Images/HR_vs_Downsample_interpolation_ bilinear.png}}
    \caption{Comparison between the original high-resolution (HR) image and its bilinearly downsampled counterpart. 
    Bilinear interpolation produces a smooth but visibly blurred low-resolution image, illustrating the information loss that the super-resolution network aims to reconstruct.}
    \label{fig:hr-vs-downsample-bilinear}
\end{figure}

\paragraph{ResizeByScale.}
The \texttt{ResizeByScale} layer replaces conventional pooling by resizing feature maps according to a continuous scale factor using bilinear interpolation with optional anti-aliasing. 
This approach preserves spatial alignment and enables flexible downsampling without strided convolutions or max-pooling.

\paragraph{ResizeToMatch.}
The \texttt{ResizeToMatch} layer adjusts one tensor to match the spatial size of another before concatenation in skip connections. 
It ensures consistent feature alignment between encoder and decoder stages, allowing smooth reconstruction across varying scales.

\paragraph{ClipAdd.}
In the reconstruction stage, the \texttt{ClipAdd} layer adds the predicted residual to the low-resolution input and the clips values to the valid intensity range \([0,1]\). 
This operation implements residual learning in a stable and numerically safe form.

Together, these layers replace non-differentiable resizing operations with smooth, learnable interpolation, 
enabling the network to adapt its depth and spatial scaling continuously during training.

%-----------------------------------------------------------

\section{Loss Functions and Evaluation Metrics}
\label{sec:prelim_loss_metrics}

Learning objectives and evaluation criteria are defined through task-specific loss functions and metrics that quantify model performance.  
Super-resolution (SR) focuses on pixel fidelity and perceptual similarity, whereas segmentation emphasizes spatial overlap and classification accuracy.

\subsection{Super-Resolution Objectives}

\textbf{Loss Functions.}
A loss function \( L(\hat{Y}, Y) \) measures the discrepancy between the predicted image \(\hat{Y}\) and the ground truth \(Y\).
Common choices include:

\begin{itemize}
    \item \textbf{L1 Loss (Mean Absolute Error):}
    \[
    L_{L1} = \frac{1}{N} \sum_{i=1}^{N} |\hat{Y}_i - Y_i|
    \]
    Encourages sharp reconstructions and is less sensitive to outliers than L2.

    \item \textbf{L2 Loss (Mean Squared Error):}
    \[
    L_{L2} = \frac{1}{N} \sum_{i=1}^{N} (\hat{Y}_i - Y_i)^2
    \]
    Penalizes large deviations more strongly, leading to smoother but sometimes overly averaged outputs.

    \item \textbf{Charbonnier Loss (Smooth L1):}
    \[
    L_{Charbonnier} = \frac{1}{N} \sum_{i=1}^{N} \sqrt{(\hat{Y}_i - Y_i)^2 + \epsilon^2}
    \]
    A differentiable L1 variant combining robustness to outliers with stable gradients near convergence \cite{charbonnier1994two}.
\end{itemize}

\textbf{Evaluation Metrics.}
Performance is quantified by reconstruction fidelity and perceptual similarity:

\begin{itemize}
    \item \textbf{Peak Signal-to-Noise Ratio (PSNR):}
    \[
    \text{PSNR} = 10 \cdot \log_{10}\!\left(\frac{MAX_I^2}{\text{MSE}}\right)
    \]
    Higher PSNR indicates closer pixel-level correspondence.

    \item \textbf{Structural Similarity Index (SSIM):}
    Evaluates luminance, contrast, and structural similarity between reference and reconstructed images; ranges from \(-1\) to \(1\) \cite{wang2004ssim}.

    \item \textbf{Multi-Scale SSIM (MS-SSIM):}
    Extends SSIM across multiple resolutions for perceptual robustness \cite{wang2003multiscale}.
\end{itemize}

\subsection{Segmentation Objectives}

\textbf{Loss Functions.}
Segmentation losses encourage accurate pixel classification and region overlap:

\begin{itemize}
    \item \textbf{Binary Cross-Entropy (BCE):}
    \[
    L_{BCE} = -\frac{1}{N}\sum_{i=1}^{N} \big[y_i \log(\hat{y}_i) + (1 - y_i)\log(1 - \hat{y}_i)\big]
    \]
    Measures per-pixel classification error.

    \item \textbf{Dice Loss:}
    \[
    L_{Dice} = 1 - \frac{2 \sum_i y_i \hat{y}_i + \epsilon}{\sum_i y_i + \sum_i \hat{y}_i + \epsilon}
    \]
    Maximizes overlap between prediction and ground truth; robust to class imbalance \cite{dice1945measures}.
\end{itemize}

\textbf{Evaluation Metrics.}
Quantitative accuracy is measured using region overlap metrics:

\begin{itemize}
    \item \textbf{Dice Coefficient:}
    \[
    \text{Dice} = \frac{2 \times |P \cap G|}{|P| + |G|}
    \]
    Measures similarity between predicted and true masks \cite{dice1945measures}.

    \item \textbf{Intersection-over-Union (IoU):}
    \[
    \text{IoU} = \frac{|P \cap G|}{|P \cup G|}
    \]
    Standard in segmentation benchmarks such as PASCAL VOC \cite{everingham2010pascal}.
\end{itemize}


%-----------------------------------------------------------
\section{Summary}
\label{sec:prelim_summary}

This chapter presented the theoretical background required to understand 
the adaptive-depth U-Net methodology. 
It covered the fundamental principles of super-resolution and segmentation, 
explained the core deep learning mechanisms behind the U-Net, 
and described the loss functions and evaluation metrics used to assess model performance.